\chapter{Channel Pulse Response and Modeling}

\section{Overview}

This chapter develops time-domain channel models and examines
the impulse and pulse responses of communication channels

\section{Channel LTI Model}
Assume input discrete-time sequence is \(x[n]\), corresponding continuous-time signal \(x(t)\) could be written as

\begin{equation}
x(t)=\sum_{n=-\infty}^{\infty}x[n]\,w_T(t-nT)
\end{equation}

where \(w_T(t)\) is a \textbf{window function} with unit amplitude and duration \(T\), corresponding to one symbol period.

\begin{equation}
w_T(t)=u(t)-u(t-T).
\end{equation}

The received signal is

\begin{equation}
y(t)=x(t)*h_{\mathrm{ch}}(t)=\sum_{n=-\infty}^{\infty}x[n]\,w_T(t-nT)*h_{\mathrm{ch}}(t)
\end{equation}

Define the \textbf{Channel Pulse Response} as

\begin{equation}
p_{\mathrm{ch}}(t)
=
w_T(t)*h_{\mathrm{ch}}(t).
\end{equation}

Then

\begin{equation}
y(t)
=
\sum_{n=-\infty}^{\infty}
x[n]\,
p_{\mathrm{ch}}(t-nT).
\end{equation}

Let \(t_0\) be the time at which \(p_{\mathrm{ch}}(t)\) reaches its peak, and define

\begin{equation}
p_{\mathrm{ch}}[n]
=
p_{\mathrm{ch}}(t_0+nT).
\end{equation}

Thus,

\begin{equation}
p_{\mathrm{ch}}[0]
=
p_{\mathrm{ch}}(t_0)
\end{equation}

is the main cursor.

Sampling the received signal at

\begin{equation}
t=t_0+kT,
\end{equation}

gives

\begin{equation}
y[k]
=
y(t_0+kT).
\end{equation}

Therefore,

\begin{equation}
y[k]
=
\sum_{n=-\infty}^{\infty}
x[n]\,
p_{\mathrm{ch}}(t_0+(k-n)T).
\end{equation}

Using the definition of \(p_{\mathrm{ch}}[n]\),

\begin{equation}
y[k]
=
\sum_{n=-\infty}^{\infty}
x[n]\,
p_{\mathrm{ch}}[k-n].
\end{equation}

Hence,

\begin{equation}
y[k]
=
x[n]*p_{\mathrm{ch}}[n].
\end{equation}


% ---------------------------------------------------------
% Bibliography
% ---------------------------------------------------------

